\documentclass[11pt,letterpaper]{article}
\usepackage[T1]{fontenc}
\usepackage{lmodern}
\usepackage[margin=1in,headheight=14pt]{geometry}
\usepackage{amsmath,amssymb,amsthm,mathtools,bm,mathrsfs}
\usepackage{microtype,booktabs,longtable,array,enumitem}
\usepackage[dvipsnames]{xcolor}
\usepackage{fancyhdr}
\usepackage[breaklinks=true,colorlinks=true,linkcolor=MidnightBlue,citecolor=MidnightBlue,urlcolor=MidnightBlue]{hyperref}
\usepackage{bookmark}
\hypersetup{pdftitle={Polynomial Lerch Selectors: Endpoint Obstructions, Resonant Spectral Collapse, and Stieltjes Confluence},pdfauthor={Research continuation for the ProveIt project}}
\pagestyle{fancy}\fancyhf{}
\fancyhead[L]{\small Polynomial Lerch selectors}
\fancyhead[R]{\small ProveIt research continuation}
\fancyfoot[C]{\thepage}
\setlength{\emergencystretch}{2em}
\setlength{\parskip}{4pt}
\numberwithin{equation}{section}
\newtheorem{theorem}{Theorem}[section]
\newtheorem{lemma}[theorem]{Lemma}
\newtheorem{proposition}[theorem]{Proposition}
\newtheorem{corollary}[theorem]{Corollary}
\theoremstyle{definition}
\newtheorem{definition}[theorem]{Definition}
\newtheorem{example}[theorem]{Example}
\theoremstyle{remark}\newtheorem{remark}[theorem]{Remark}
\newcommand{\C}{\mathbb C}\newcommand{\Q}{\mathbb Q}
\newcommand{\dd}{\,\mathrm d}
\newcommand{\Li}{\operatorname{Li}}
\newcommand{\Res}{\operatorname{Res}}
\newcommand{\supp}{\operatorname{supp}}
\newcommand{\T}{\mathcal T}
\newcommand{\E}{\mathcal E}
\newcommand{\HFE}{\mathsf H}
\newcommand{\src}[1]{\texttt{\detokenize{#1}}}
\DeclareMathOperator{\Span}{span}
\title{\textbf{Polynomial Lerch Selectors}\\[7pt]
\large Endpoint obstructions, resonant spectral collapse,\\
and Stieltjes confluence}
\author{Research continuation prepared for\\Vladimir Reshetnikov's ProveIt project}
\date{11 October 2026}
\begin{document}
\maketitle
\begin{abstract}
We answer the polynomial-selector question posed in the incoming
\emph{Affine Gamma Polarization} report for every positive integer
interval length and every nonzero twist in the closed unit disk.
A finite vector of paired endpoint derivatives is exactly the support
of the shifted Lerch functions in the transform, modulo a specified
finite exponential remainder. Off resonance, the unique single-order
weights are classical Frobenius--Euler polynomials; at resonance they
are Bernoulli polynomials, with an arbitrary constant. Functional
independence over $\C(s)$ makes this a classification of analytic
functions, not merely of formal coordinates.

Eulerian polynomials give a sharp endpoint obstruction: off resonance
a nonzero single-order weight cannot vanish to order two at both
endpoints. At resonance, double zeros occur precisely for the even
Bernoulli weights of degree at least four, and triple zeros never occur.
An explicit finite mixture replaces the impossible scalar confluence;
its number of spectral coordinates drops exactly from $r+1$ to one.
We derive every removable-integer spectral jet through finite
generalized-harmonic coefficients and obtain Euler-weighted Gamma
integrals with one odd zeta value. A joint holomorphic renormalization
at the Hurwitz pole supplies every mixed twist/Stieltjes jet and
anchored shift primitives.

Classical polynomial bases and the classical Lerch expansion are
credited explicitly. The contribution is their exact selector
classification, endpoint and confluence consequences, and the resulting
identity calculus in the stated source programme. We also solve an
explicit coefficient problem in arXiv:2609.13796v1 and correct three
normalization/sign statements relative to its defining series. The package includes
969 finite exact assertions and 33 high-precision diagnostics. These
are not interval certificates or proof-assistant verification. No global
priority claim or resolution of the $S_6$, revised $S_8$, or general
unequal-tail prefix problems is made.
\end{abstract}
\newpage
\tableofcontents
\newpage

\section{The source question and the scope of the answer}
\label{pls:sec:scope}

The canonical ProveIt manuscript separates its proven identities,
experimental candidates, and pending analytic integrations
\cite{Manuscript}. The incoming \emph{Affine Gamma Polarization} report
extends the Gamma square sector to weighted Hurwitz and Lerch
transforms. Its Section~11, questions~4 and~5, asks for a functional
classification of polynomial weights on integer intervals that leave
one shifted zeta or polylogarithm order plus an elementary remainder,
and asks how the twisted answer approaches the Bernoulli selectors
\cite{Affine}. These are the specific continuation questions addressed
here. We do not repeat that report's Gamma normal-form theorem.

The phrase ``one order'' needs a coefficient field. At a single
numerical value of $s$, identities between constants could obscure the
issue. With arbitrary meromorphic coefficients every nonzero function
could trivially be called a selector. We instead allow coefficients in
$\C(s)$ and fix the finite exponential remainder arising from integer
translation. This is the same functional, rather than numerical,
standard used in the predecessor's Bernoulli classification.

Let $a>0$, $L\ge1$ be an integer, $P\in\C[u]$, and
$0<|z|\le1$. The central transform is
\begin{equation}\label{pls:eq:transform}
 \T_P(z,s;a,L)=\int_0^L P(u)\Phi(z,s,a+u)\dd u,
 \qquad
 \Phi(z,s,x)=\sum_{n=0}^\infty\frac{z^n}{(n+x)^s}.
\end{equation}
For $|z|<1$ the series is entire in $s$. At $|z|=1$ the identities
are first proved for $\Re s>1$ and then continued. For $z\ne1$ the
continued transform is entire; for $z=1$ it is a weighted Hurwitz-zeta
integral with only a possible pole at $s=1$. Positive bases use their
real logarithms. No branch ambiguity in $(n+x)^{-s}$ is introduced.

Theorem~\ref{pls:thm:classification} solves the exact selector problem.
Theorem~\ref{pls:thm:flatness} proves the endpoint obstruction and its
resonant exception. Theorem~\ref{pls:thm:mixing} describes constructive
confluence with an exact support count.
Theorems~\ref{pls:thm:jets} and~\ref{pls:thm:confluence} give the
spectral and Stieltjes calculi. The concrete Gamma and quartic-weight
identities in Sections~\ref{pls:sec:gamma} and~\ref{pls:sec:stieltjes}
are useful normalization tests and applications.

\paragraph{Attribution and novelty scope.}
Frobenius--Euler polynomials, their Appell property and their endpoint
basis formula are classical; see Kim and Kim \cite{KimKim}. Eulerian
real-rootedness is classical as well; modern interlacing proofs include
Chow \cite{Chow}. Hurwitz differentiation, Lerch translation, and the
logarithmic Lerch expansion are classical \cite{DLMFHurwitz,DLMFLerch}.
The predecessor already proves a single-zeta Bernoulli basis on $[1,3]$
and weighted Stieltjes formulas. We give self-contained proofs where
short, but do not claim these ingredients as discoveries. The results
below resolve the stated source questions in an explicit domain and
exhibit consequences not established in the inspected source sections.
A finite literature search cannot establish global priority.

\section{Paired endpoint derivatives are the exact spectrum}
\label{pls:sec:endpoints}

Write
\begin{equation}\label{pls:eq:F}
 F_j(s)=\prod_{k=1}^{j+1}(s-k),\qquad j\ge0,
 \qquad \lambda=z^L.
\end{equation}
For a polynomial of degree $D$, set
\begin{equation}\label{pls:eq:c}
 c_j(P;\lambda,L)=P^{(j)}(L)-\lambda P^{(j)}(0),
 \qquad 0\le j\le D.
\end{equation}
The elementary remainder space in the next results is
\begin{equation}\label{pls:eq:elementary}
 \E_{a,L}=\Span_{\C(s)}\{(a+\ell)^{-s}:0\le\ell<L\}.
\end{equation}
The twist is fixed while taking this span. Coefficients may depend on
that fixed twist and on the fixed interval parameters.

\begin{proposition}[Finite endpoint transform]\label{pls:prop:endpoint}
As meromorphic functions of $s$,
\begin{align}
 \T_P(z,s;a,L)
 &=-\sum_{j=0}^D\frac{P^{(j)}(L)\Phi(z,s-j-1,a+L)
                  -P^{(j)}(0)\Phi(z,s-j-1,a)}{F_j(s)}
                  \label{pls:eq:endpoint}\\
 &=-\lambda^{-1}\sum_{j=0}^D
        \frac{c_j(P;\lambda,L)}{F_j(s)}\Phi(z,s-j-1,a)
      +\mathcal R_P(z,s;a,L),\label{pls:eq:spectrum}
\end{align}
where the explicitly elementary remainder is
\begin{equation}\label{pls:eq:remainder}
 \mathcal R_P=\lambda^{-1}\sum_{j=0}^D
     \frac{P^{(j)}(L)}{F_j(s)}
          \sum_{\ell=0}^{L-1}z^\ell(a+\ell)^{j+1-s}.
\end{equation}
All apparent poles other than the possible Hurwitz pole at $s=1$
cancel in the complete expression.
\end{proposition}
\begin{proof}
Termwise differentiation in an initial convergence domain gives
$\partial_x\Phi(z,s,x)=-s\Phi(z,s+1,x)$. Successive integration by
parts terminates after $D+1$ steps. At each step the antiderivative
contributes the next factor $-(s-k)^{-1}$ while integration by parts
contributes another minus sign to the continuing integral. Thus the
boundary sum has the uniform minus sign in
\eqref{pls:eq:endpoint}. Equivalently,
$(-1)^j/(1-s)_{j+1}=-1/F_j(s)$.

Integer translation of the absolutely convergent series gives
\[
 \Phi(z,\sigma,a)=z^L\Phi(z,\sigma,a+L)
                  +\sum_{\ell=0}^{L-1}z^\ell(a+\ell)^{-\sigma}.
\]
Substitution proves \eqref{pls:eq:spectrum}. Analytic continuation
then proves the formulas everywhere they are asserted. For $z\ne1$
on the unit circle, the Mellin representation with denominator
$1-ze^{-t}$ continues the function to an entire function of $s$:
near zero its numerator is analytic, and the zeros of $1/\Gamma(s)$
cancel the Mellin poles at nonpositive integers. The corresponding
assertion at $z=1$ is the usual Hurwitz continuation. Integration over
a compact positive interval preserves these properties.
\end{proof}

\begin{lemma}[Functional independence]\label{pls:lem:independence}
For distinct nonnegative integers $j_1,\ldots,j_m$, the functions
\[
 \Phi(z,s-j_1-1,a),\ldots,\Phi(z,s-j_m-1,a),
 \quad (a+\ell)^{-s}\ (0\le\ell<L)
\]
are linearly independent over $\C(s)$.
\end{lemma}
\begin{proof}
Clear rational denominators in a putative relation. For real $s$
sufficiently large, expand it into a Dirichlet series
\[
 \sum_{n\ge0}Q_n(s)(a+n)^{-s}=0.
\]
The polynomials $Q_n(s)$ have uniformly bounded degree in $s$. For
$n\ge L$ they have the form
\[
 Q_n(s)=z^n\sum_{\nu=1}^m p_\nu(s)(a+n)^{j_\nu+1},
\]
where the $p_\nu$ are fixed polynomials. The finitely many smaller
indices contain the elementary coefficients as well. There are
constants $C,M,K$ for which
$|Q_n(s)|\le C(1+s)^M(1+n)^K$ on the positive real axis.

If one coefficient polynomial is nonzero, choose the least such
index $n_0$. Divide the identity by $(a+n_0)^{-s}$. The tail tends
to zero exponentially, up to a fixed polynomial in $s$. Indeed, for
large $s$ one can split the exponent in half and bound it by
\[
 C'(1+s)^M
 \left(\frac{a+n_0}{a+n_0+1}\right)^{s/2}
 \sum_{n>n_0}(1+n)^K
          \left(\frac{a+n_0}{a+n}\right)^{s/2},
\]
where the final sum is uniformly bounded once $s/2>K+2$.
A nonzero polynomial $Q_{n_0}(s)$ cannot decay exponentially.
Thus every $Q_n$ is zero.

For infinitely many $n\ge L$, division by $z^n\ne0$ gives
$\sum_\nu p_\nu(s)(a+n)^{j_\nu+1}=0$. A polynomial in $a+n$
with infinitely many roots is identically zero; hence each $p_\nu=0$.
The remaining finite exponential coefficients are then zero too.
\end{proof}

\begin{corollary}[No hidden spectral cancellation]\label{pls:cor:support}
Modulo $\E_{a,L}$, the transform has precisely the shifted orders
$s-j-1$ for which $c_j(P;\lambda,L)\ne0$. This support is unique.
\end{corollary}

For positive integer $a=A$, the relation
\begin{equation}\label{pls:eq:integerLi}
 \Phi(z,\sigma,A)=z^{-A}\left(\Li_\sigma(z)
                -\sum_{n=1}^{A-1}\frac{z^n}{n^\sigma}\right)
\end{equation}
turns this into an exact polylogarithmic statement. The elementary
space is then enlarged, explicitly, to
$\Span_{\C(s)}\{n^{-s}:1\le n<A+L\}$.
The same least-index argument proves independence in that form.

\section{Classification off and at resonance}
\label{pls:sec:classification}

We call $\lambda=z^L=1$ \emph{resonance}. This concerns the interval
translation, not only the point $z=1$. Nontrivial roots of unity can
therefore be resonant, while their Lerch functions remain entire in $s$.

Define the classical Frobenius--Euler polynomials by
\begin{equation}\label{pls:eq:FEgen}
 \frac{(1-\lambda)e^{xt}}{e^t-\lambda}
 =\sum_{d=0}^\infty\HFE_d(x;\lambda)\frac{t^d}{d!},
 \qquad \lambda\ne1.
\end{equation}
They are monic and satisfy
\begin{align}
 \partial_x\HFE_d(x;\lambda)&=d\HFE_{d-1}(x;\lambda),
                 \label{pls:eq:Appell}\\
 \HFE_d(1;\lambda)-\lambda\HFE_d(0;\lambda)
                  &=(1-\lambda)\mathbf1_{d=0}.
                 \label{pls:eq:FEboundary}
\end{align}
These identities follow immediately from the generating function.
The notation $\HFE_d$ is distinguished from generalized harmonic
numbers $H_n^{(r)}$ used later.

\begin{theorem}[All single-order polynomial selectors]
\label{pls:thm:classification}
Fix $a>0$, integer $L\ge1$, and $0<|z|\le1$. Let $N\ge1$.
A polynomial $P$ has a nonzero single-order representation
\begin{equation}\label{pls:eq:singledef}
 \T_P(z,s;a,L)=A(s)\Phi(z,s-N,a)+E(s),
 \qquad 0\ne A\in\C(s),\quad E\in\E_{a,L},
\end{equation}
if and only if one of the following holds.

If $\lambda\ne1$, then
\begin{equation}\label{pls:eq:nonresclass}
 P(u)=C L^{N-1}\HFE_{N-1}(u/L;\lambda),\qquad C\ne0.
\end{equation}
In particular $\deg P=N-1$, and
$A(s)=-\lambda^{-1}(1-\lambda)C(N-1)!/F_{N-1}(s)$.
No nonzero polynomial has an elementary-only transform.

If $\lambda=1$, then
\begin{equation}\label{pls:eq:resclass}
 P(u)=C L^N B_N(u/L)+C_0,\qquad C\ne0,
\end{equation}
where $B_N$ is the Bernoulli polynomial. Now $\deg P=N$ and
$A(s)=-CLN!/F_{N-1}(s)$. The elementary-only polynomials are
exactly the constants.
\end{theorem}
\begin{proof}
For $\lambda\ne1$, \eqref{pls:eq:Appell}--\eqref{pls:eq:FEboundary}
show that the polynomial $L^d\HFE_d(u/L;\lambda)$ has only one
nonzero endpoint coordinate, namely $c_d=(1-\lambda)d!$.
Since these monic polynomials form a basis, every $P$ of degree at
most $D$ has the classical endpoint expansion
\begin{equation}\label{pls:eq:FEbasis}
 P(u)=\sum_{d=0}^D
       \frac{c_d(P;\lambda,L)L^d}{(1-\lambda)d!}
                         \HFE_d(u/L;\lambda).
\end{equation}
Corollary~\ref{pls:cor:support} proves that a single order is
possible exactly when one coefficient in this expansion survives.
The highest endpoint coordinate of a nonzero degree-$D$ polynomial
is $(1-\lambda)P^{(D)}$, so its sole coordinate must have $d=D$.
This proves the nonresonant assertion.

At resonance use
$B_n(1)-B_n(0)=\mathbf1_{n=1}$, which follows from
$te^{xt}/(e^t-1)$. It gives the anchored expansion
\begin{equation}\label{pls:eq:Bbasis}
 P(u)=P(0)+\sum_{N=1}^D
 \frac{c_{N-1}(P;1,L)L^{N-1}}{N!}
                      \bigl(B_N(u/L)-B_N(0)\bigr).
\end{equation}
Only the coordinate $j=N-1$ is nonzero for a nonconstant basis term.
The kernel consists of constants: alternatively, a polynomial with
all paired derivatives zero must be constant, by looking successively
at its highest nonconstant degree. The independence lemma again
turns this basis statement into the functional classification.
\end{proof}

\begin{remark}[A finite interpolation problem, with a singular limit]
On the monomial basis $1,u,\ldots,u^D$, the endpoint map
$P\mapsto(c_0,\ldots,c_D)$ is upper triangular and has determinant
\begin{equation}\label{pls:eq:endpointdet}
 (1-\lambda)^{D+1}\prod_{j=0}^D j!.
\end{equation}
At $\lambda=1$ it has rank $D$, with constants as kernel. The minor
with rows $c_0,\ldots,c_{D-1}$ and columns $u,\ldots,u^D$ has
\begin{equation}\label{pls:eq:resdet}
 \det=L^D\prod_{j=1}^D j!\ne0.
\end{equation}
Thus the degree shift in Theorem~\ref{pls:thm:classification} is an
exact change of rank, not a numerical instability.
\end{remark}

\section{Endpoint zeros and the Eulerian obstruction}
\label{pls:sec:flatness}

The source's Gamma-derived weights vanish with their first derivatives
at both endpoints. Such a condition is not innocuous under twisting.
We now classify precisely when it is compatible with one order.

Define Eulerian polynomials $A_d$ for $d\ge1$ by
\begin{equation}\label{pls:eq:Eulerian}
 \sum_{n=1}^\infty n^d\lambda^n
       =\frac{\lambda A_d(\lambda)}{(1-\lambda)^{d+1}},
 \qquad |\lambda|<1.
\end{equation}
Here $A_1=1$, $A_2=1+\lambda$, and
$A_3=1+4\lambda+\lambda^2$. Differentiating the generating series
with $\lambda\partial_\lambda$ gives
\begin{equation}\label{pls:eq:Eulerianrec}
 A_d=(1+(d-1)\lambda)A_{d-1}
          +\lambda(1-\lambda)A_{d-1}'.
\end{equation}

\begin{lemma}[The endpoint Eulerian polynomial]\label{pls:lem:Eulerian}
For $d\ge1$ and $\lambda\ne1$,
\begin{equation}\label{pls:eq:FEEulerian}
 \HFE_d(0;\lambda)=\frac{(-1)^d A_d(\lambda)}{(1-\lambda)^d}.
\end{equation}
Every zero of $A_d$ is simple and negative. Consecutive $A_d,A_{d-1}$, for $d\ge2$,
have no common zero.
\end{lemma}
\begin{proof}
For $|\lambda|<1$, write the left side of \eqref{pls:eq:FEgen} at
$x=0$ as
$(1-\lambda)\sum_{n\ge0}\lambda^n e^{-(n+1)t}$.
Taking its $d$th derivative at zero and using
\eqref{pls:eq:Eulerian} proves \eqref{pls:eq:FEEulerian}.
Rational continuation proves it for every $\lambda\ne1$.

For completeness, real-rootedness follows directly by induction from
\eqref{pls:eq:Eulerianrec}. The polynomial $A_{d-1}$ is monic of
degree $d-2$, has value one at zero, and, inductively, has simple
negative roots $r_1<\cdots<r_{d-2}$. At each such root,
\[
 A_d(r_i)=r_i(1-r_i)A_{d-1}'(r_i)\ne0.
\]
These values alternate in sign. There is one zero of $A_d$ in each
successive gap, one between $r_{d-2}$ and zero, and one to the left
of $r_1$, by comparing its monic leading sign at $-\infty$.
These $d-1$ distinct negative zeros exhaust its degree. The base
$A_2=1+\lambda$ starts the argument; $A_1$ has no roots. The same
displayed identity proves the asserted absence of common zeros.
\end{proof}

\begin{theorem}[Sharp endpoint multiplicity obstruction]
\label{pls:thm:flatness}
A nonzero polynomial single-order selector has the following endpoint
possibilities.

For $\lambda\ne1$, its degree-$d$ weight can vanish at both endpoints
only if $d\ge2$ and $A_d(\lambda)=0$. These zeros of the weight are
then simple at both endpoints. In particular no nonzero single-order
selector satisfies
\begin{equation}\label{pls:eq:doubleflat}
 P(0)=P'(0)=P(L)=P'(L)=0.
\end{equation}
For real $0<z<1$, even simultaneous endpoint values zero are impossible.

At $\lambda=1$, a single-order selector can satisfy
\eqref{pls:eq:doubleflat} if and only if it is a nonzero multiple of
\begin{equation}\label{pls:eq:Bernflat}
 B_{2r}(u/L)-B_{2r},\qquad r\ge2.
\end{equation}
Both endpoint zeros then have order exactly two. No nonzero
single-order selector in either regime has zeros of order three at
both endpoints.
\end{theorem}
\begin{proof}
For $d\ge1$, \eqref{pls:eq:FEboundary} makes the two endpoint values
proportional by the nonzero factor $\lambda$. Their vanishing is
therefore equivalent to $A_d(\lambda)=0$. The cases $d=0,1$ have
no vanishing endpoint value. If $d\ge2$ and the value is zero, the
Appell identity makes the derivative proportional to
$\HFE_{d-1}(0;\lambda)$, which is nonzero by the lemma.
The derivative at the other endpoint is again multiplied by
$\lambda$, so it is also nonzero. If $0<z<1$, then $\lambda>0$,
where $A_d$ has no root.

At resonance the classification permits only
$C L^N B_N(u/L)+C_0$. For $N=1$ its two endpoint values differ,
so they cannot both be zero. For $N\ge2$, requiring them both to
vanish fixes the constant and gives a multiple of
$B_N(u/L)-B_N$. Its derivative vanishes at both endpoints exactly
when $B_{N-1}(0)=B_{N-1}(1)=0$, namely when $N\ge4$ is even.
Its second derivative is a nonzero multiple of $B_{N-2}$, a nonzero
even Bernoulli number. The zeros have exactly order two.
\end{proof}

\begin{example}[A nontrivial algebraic twist]\label{pls:ex:algebraic}
For $L=1$, $z=\sqrt3-2$, and $d=3$, the identity
$A_3(z)=0$ gives the exact weight
\begin{equation}\label{pls:eq:algebraicP}
 \HFE_3(u;\sqrt3-2)
 =u(u-1)\left(u-\frac{1+\sqrt3}{2}\right).
\end{equation}
It selects only $\Li_{s-4}(\sqrt3-2)$ modulo the elementary
remainder. Both endpoint zeros are simple. This distinguishes
``endpoint vanishing'' from the stronger and impossible off-resonant
condition \eqref{pls:eq:doubleflat}.
\end{example}

\section{Constructive confluence and exact spectral collapse}
\label{pls:sec:mixing}

The Frobenius--Euler family does not approach the Bernoulli family by
an innocuous scalar normalization. First we make that obstruction
precise, and then give an exact finite replacement.

\begin{proposition}[No nonconstant scalar limit]\label{pls:prop:scalar}
Fix $d\ge1$. Suppose that
$C(\lambda)\HFE_d(u;\lambda)$ has a finite coefficientwise limit
as $\lambda\to1$. That limit is constant. The statement also holds
after the fixed scaling $u\mapsto u/L$.
\end{proposition}
\begin{proof}
The recurrence \eqref{pls:eq:Eulerianrec} gives $A_d(1)=d!$.
Consequently
\[
 \HFE_d(0;\lambda)\sim\frac{(-1)^d d!}{(1-\lambda)^d}.
\]
Boundedness of the constant coefficient forces
$C(\lambda)=O(|1-\lambda|^d)$. By the Appell property, the
coefficient of $u^j$ is
$\binom dj\HFE_{d-j}(0;\lambda)$, with pole order at most $d-j$.
After multiplication by $C(\lambda)$ every coefficient with $j\ge1$
tends to zero. The constant coefficient may have a finite limit.
\end{proof}

This proposition concerns one fixed degree and scalar multiplication
only. It does not forbid finite combinations of different degrees,
subtraction of lower-degree counterterms, or unbounded-degree limits.
The finite combination below is especially sparse.

\begin{theorem}[Minimal mixing for the fixed even Bernoulli weight]
\label{pls:thm:mixing}
Let $N=2r\ge4$ and $W_N(u)=B_N(u)-B_N$. For $z\ne1$,
\begin{equation}\label{pls:eq:mixing}
 \boxed{\displaystyle
 W_N(u)=\sum_{\substack{2\le j\le N\\j\ {\rm even}}}
       \binom Nj B_{N-j}\HFE_j(u;z)
   +\frac{N(1+z)}{2(1-z)}\HFE_{N-1}(u;z).}
\end{equation}
For the transform on $[0,1]$, the number of distinct shifted Lerch
orders modulo the elementary remainder is exactly
\begin{equation}\label{pls:eq:supportcount}
 \begin{cases}
 r+1,& z\notin\{0,1,-1\},\quad |z|\le1,\\
 r,& z=-1,\\
 1,& z=1.
 \end{cases}
\end{equation}
The mixture in \eqref{pls:eq:mixing} is unique in the
Frobenius--Euler basis, and no representation of this fixed transform
can use fewer distinct shifted orders in the declared coefficient
system. On $[0,L]$ with $W_N(u/L)$ the same count holds with $z$
replaced by $\lambda=z^L$.
\end{theorem}
\begin{proof}
Put $b_j=W_N^{(j)}(1)$. The even nonzero endpoint derivatives are
\begin{equation}\label{pls:eq:bj}
 b_j=\frac{N!}{(N-j)!}B_{N-j},\qquad j=2,4,\ldots,N.
\end{equation}
There is also the odd derivative $b_{N-1}=N!/2$; every other odd
derivative is zero, as are $b_0,b_1$.
For even $j$, the left and right endpoint derivatives coincide.
For $j=N-1$ they are opposite. Thus
\begin{equation}\label{pls:eq:mixc}
 c_j=(1-z)b_j\quad(j\text{ even}),\qquad
 c_{N-1}=(1+z)N!/2,
\end{equation}
and all other coordinates vanish. Substitute these coordinates into
\eqref{pls:eq:FEbasis} to obtain \eqref{pls:eq:mixing}.
Every even Bernoulli number in \eqref{pls:eq:bj}, including $B_0$,
is nonzero. The independence lemma therefore gives the exact count
\eqref{pls:eq:supportcount}, including its specializations at $z=\pm1$.
Scaling the interval multiplies $b_j$ by $L^{-j}$ and changes the
paired endpoint factor to $z^L$, without changing the count.
\end{proof}

For example,
\begin{equation}\label{pls:eq:quarticmix}
 u^2(1-u)^2=\HFE_4(u;z)+\HFE_2(u;z)
                   +\frac{2(1+z)}{1-z}\HFE_3(u;z).
\end{equation}
At $z=-1$ this is simply $E_4(u)+E_2(u)$, where $E_d$ is the
ordinary Euler polynomial. As $z\to1$, the polynomial on the left
stays fixed while the individual terms on the right need not have
finite limits. This is the precise cancellation missed by a scalar
confluence argument.

For $a=1$, an explicit version of the transform is
\begin{align}
 \int_0^1W_N(u)\Phi(z,s,1+u)\dd u
 ={}&-\frac{N!(1+z)}{2z^2F_{N-1}(s)}\Li_{s-N}(z)
       \notag\\
 &-\frac{1-z}{z^2}
       \sum_{\substack{2\le j\le N\\j\ {\rm even}}}
                      \frac{b_j\Li_{s-j-1}(z)}{F_j(s)}
   +\frac1z\sum_{j=2}^N\frac{b_j}{F_j(s)}.
\label{pls:eq:BernLi}
\end{align}
It follows directly from \eqref{pls:eq:spectrum}. Its limit at
$z=1$ is initially justified for $\Re s>1$. At $s=1$ the limiting
Hurwitz pole requires the compensation in
Section~\ref{pls:sec:confluence}; it must not be suppressed by
substituting into individual summands.

\section{Removable integer orders and harmonic jet coefficients}
\label{pls:sec:jets}

The endpoint denominators vanish at positive integers, even when the
integral is entire in $s$. The selector classification allows these
limits to be evaluated using derivatives at just one polylogarithm
order. We give a complete finite formula rather than leaving a
removable $0/0$ expression.

For $z\ne0,1$ in the closed unit disk define
\begin{align}
 T_d(z,s)&=\int_0^1\HFE_d(u;z)\Phi(z,s,1+u)\dd u,
                       \label{pls:eq:Td}\\
 \Pi_d(s)&=\prod_{k=1}^{d+1}(s-k),\label{pls:eq:Pid}\\
 R_d(s,z)&=\sum_{j=0}^d\HFE_d^{(j)}(1;z)
                           \prod_{k=j+2}^{d+1}(s-k).
                       \label{pls:eq:Rd}
\end{align}
The superscript $(j)$ in \eqref{pls:eq:Rd} differentiates the
polynomial variable. Empty products equal one. The endpoint formula
and $\Phi(z,\sigma,1)=z^{-1}\Li_\sigma(z)$ yield
\begin{equation}\label{pls:eq:numerator}
 \Pi_d(s)T_d(z,s)=\frac{R_d(s,z)}z
          -\frac{(1-z)d!}{z^2}\Li_{s-d-1}(z)
          =:N_d(s,z).
\end{equation}
The numerator is entire in $s$ and vanishes at each
$m=1,\ldots,d+1$.

\begin{theorem}[Every integer spectral jet of a selector]
\label{pls:thm:jets}
For $1\le m\le d+1$, put
\begin{equation}\label{pls:eq:pm}
 p_{d,m}=(-1)^{d+1-m}(m-1)!(d+1-m)!,
\end{equation}
and define the analytic germ
\begin{align}
 E_{d,m}(\varepsilon)
 &=\prod_{\substack{1\le k\le d+1\\k\ne m}}
                  \left(1+\frac{\varepsilon}{m-k}\right)^{-1}
                  \label{pls:eq:Em}\\
 &=\exp\left(\sum_{r\ge1}
 \frac{(-1)^rH_{m-1}^{(r)}+H_{d+1-m}^{(r)}}r
                                     \varepsilon^r\right).
                  \label{pls:eq:Emharmonic}
\end{align}
Here $H_n^{(r)}=\sum_{k=1}^n k^{-r}$ and $H_0^{(r)}=0$.
For every integer $h\ge0$,
\begin{equation}\label{pls:eq:alljets}
 \boxed{\displaystyle
 \partial_s^hT_d(z,m)=\frac{h!}{p_{d,m}}
 [\varepsilon^{h+1}]
          E_{d,m}(\varepsilon)N_d(m+\varepsilon,z).}
\end{equation}
In particular,
\begin{equation}\label{pls:eq:resvalue}
 T_d(z,m)=\frac{z^{-1}\partial_sR_d(m,z)
   -(1-z)d!z^{-2}\left.\partial_\nu\Li_\nu(z)
                                    \right|_{\nu=m-d-1}}
                      {p_{d,m}}.
\end{equation}
Thus an $h$th spectral derivative requires at most the first $h+1$
order derivatives at the single center $\nu=m-d-1$.
\end{theorem}
\begin{proof}
Factor the denominator at its simple zero:
\[
 \Pi_d(m+\varepsilon)
   =p_{d,m}\varepsilon E_{d,m}(\varepsilon)^{-1}.
\]
Since $N_d(m,z)=0$, division and Taylor coefficient extraction prove
\eqref{pls:eq:alljets}. Expanding
$-\log(1+\varepsilon/(m-k))$ gives
\eqref{pls:eq:Emharmonic}: positive differences contribute
$(-1)^rH_{m-1}^{(r)}$, while negative differences contribute
$H_{d+1-m}^{(r)}$. The constant term of $E_{d,m}$ is one, so
\eqref{pls:eq:resvalue} is the case $h=0$.
\end{proof}

The coefficient rule is executable without symbolic limits. If
$E_{d,m}=\sum_{n\ge0}e_n\varepsilon^n$ and
$A_r=(-1)^rH_{m-1}^{(r)}+H_{d+1-m}^{(r)}$, then
\begin{equation}\label{pls:eq:ecoef}
 e_0=1,\qquad e_n=\frac1n\sum_{r=1}^n A_r e_{n-r}.
\end{equation}
For example, $e_1=A_1$ and $e_2=(A_1^2+A_2)/2$.
The first derivative reads
\[
 \partial_sT_d(z,m)
 =\frac{\tfrac12\partial_s^2N_d(m,z)
               +A_1\partial_sN_d(m,z)}{p_{d,m}}.
\]
These are finite generalized-harmonic identities, not asymptotic
expansions in a large index.

\begin{example}[Degrees zero and one]
The first two formulas are
\begin{align}
 T_0(z,s)&=\frac{z^{-1}-(1-z)z^{-2}\Li_{s-1}(z)}{s-1},
                                  \label{pls:eq:d0}\\
 T_1(z,s)&=\frac{(1+z-zs)/(z(1-z))
                      -(1-z)z^{-2}\Li_{s-2}(z)}{(s-1)(s-2)}.
                                  \label{pls:eq:d1}
\end{align}
Here $\HFE_0=1$ and $\HFE_1(u;z)=u-(1-z)^{-1}$.
At the first resonance,
\begin{equation}\label{pls:eq:d0res}
 \int_1^2\Phi(z,1,x)\dd x
       =-\frac{1-z}{z^2}
                       \left.\partial_\nu\Li_\nu(z)\right|_{\nu=0}.
\end{equation}
The negative powers of $z$ in these reductions have removable limits
at $z=0$ when the original polynomial family itself is regular there.
The classification deliberately excludes $z=0$, where the infinite
Dirichlet tail disappears and the independence lemma no longer applies.
\end{example}

\section{Euler weights, Gamma primitives, and one odd zeta}
\label{pls:sec:gamma}

At $z=-1$, \eqref{pls:eq:FEgen} becomes
$2e^{ut}/(e^t+1)$, so $\HFE_d(u;-1)=E_d(u)$ is the ordinary
Euler polynomial. Write
\begin{equation}\label{pls:eq:G}
 G(u)=\log\Gamma(1+u/2)-\log\Gamma((1+u)/2).
\end{equation}
All arguments are positive and the logarithms are real. The classical
digamma difference gives \cite{DLMFPsi,DLMFHurwitz}
\begin{equation}\label{pls:eq:Gprime}
 G'(u)=\Phi(-1,1,1+u)
       =\frac12\left\{\psi((u+2)/2)-\psi((u+1)/2)\right\}.
\end{equation}
This relation follows either by separating the even and odd terms
in the alternating series or by Hurwitz translation followed by the
Laurent expansion at one.

\begin{theorem}[A complete Euler-weighted Gamma family]
\label{pls:thm:Gamma}
For $r\ge1$, define the rational number
\begin{equation}\label{pls:eq:qr}
 q_r=-\frac{\partial_sR_{2r}(1,-1)}{(2r)!}.
\end{equation}
Then the ordinary convergent integrals satisfy
\begin{align}
 K_{2r}&:=\int_0^1 E_{2r}(u)\Phi(-1,1,1+u)\dd u\notag\\
 &=q_r+
 \frac{(-1)^{r+1}(2^{2r+1}-1)(2r)!}{(2\pi)^{2r}}
                                        \zeta(2r+1),
                                        \label{pls:eq:oddzeta}\\
 \int_0^1 E_{2r-1}(u)G(u)\dd u&=-\frac{K_{2r}}{2r}.
                                        \label{pls:eq:Gammafamily}
\end{align}
The first offsets are
\begin{equation}\label{pls:eq:qexamples}
 q_1=-\frac12,\qquad q_2=\frac7{12},\qquad
 q_3=-\frac{53}{30},\qquad q_4=\frac{8417}{840}.
\end{equation}
\end{theorem}
\begin{proof}
Apply \eqref{pls:eq:resvalue} with $d=2r$, $m=1$, and $z=-1$.
The denominator is $(2r)!$, so
\[
 K_{2r}=q_r-2\left.\partial_\nu\Li_\nu(-1)\right|_{\nu=-2r}.
\]
The identity
$\Li_\nu(-1)=(2^{1-\nu}-1)\zeta(\nu)$, interpreted by
continuation at $\nu=1$, is classical. At a negative even integer
$\zeta(-2r)=0$, hence its derivative there is
$(2^{2r+1}-1)\zeta'(-2r)$. Differentiating the zeta functional
equation at that simple trivial zero gives \cite{DLMFInteger}
\[
 \zeta'(-2r)=\frac{(-1)^r(2r)!}{2(2\pi)^{2r}}\zeta(2r+1).
\]
This proves \eqref{pls:eq:oddzeta}. The Euler endpoints
$E_{2r}(0)=E_{2r}(1)=0$ and the Appell identity
$E_{2r}'=2rE_{2r-1}$ permit integration by parts in
\eqref{pls:eq:Gprime}, proving \eqref{pls:eq:Gammafamily} without
an omitted boundary constant. The finite polynomial
\eqref{pls:eq:Rd} gives \eqref{pls:eq:qexamples} by rational arithmetic.
\end{proof}

The simplest member is the explicit log-Gamma identity
\begin{equation}\label{pls:eq:Gammaexample}
 \boxed{\displaystyle
 \int_0^1\left(u-\frac12\right)
       \log\frac{\Gamma(1+u/2)}{\Gamma((1+u)/2)}\dd u
    =\frac14-\frac{7\zeta(3)}{4\pi^2}.}
\end{equation}
We do not attach a global novelty claim to individual low-degree
Gamma evaluations; the point is their uniform single-order derivation
and compatibility with all higher spectral jets.

\subsection{Every positive polygamma order}

For integer $m\ge2$, separation into even and odd terms gives
\begin{align}
 \Phi(-1,m,1+u)
 &=2^{-m}\left\{\zeta\left(m,\frac{1+u}{2}\right)
                 -\zeta\left(m,\frac{2+u}{2}\right)\right\}\notag\\
 &=\frac{(-1)^m}{2^m(m-1)!}
   \left\{\psi^{(m-1)}\left(\frac{1+u}{2}\right)
         -\psi^{(m-1)}\left(\frac{2+u}{2}\right)\right\}.
\label{pls:eq:polygamma}
\end{align}
For $1\le m\le d+1$, Theorem~\ref{pls:thm:jets} evaluates the
Euler-weighted integral through the one order derivative at $m-d-1$.
For $m\ge d+2$, there is no vanishing denominator and
\begin{equation}\label{pls:eq:eta}
 \int_0^1E_d(u)\Phi(-1,m,1+u)\dd u
       =\frac{-R_d(m,-1)+2d!\,\eta(m-d-1)}{\Pi_d(m)},
\end{equation}
where $\eta(\nu)=(1-2^{1-\nu})\zeta(\nu)$ and
$\eta(1)=\log2$. Thus every positive polygamma order has a finite
one-center evaluation, including the exceptional orders.

For $E_2(u)=u^2-u$ and $R_2(s,-1)=s-1$, the first higher-order
members are
\begin{align}
 \int_0^1(u^2-u)\Phi(-1,2,1+u)\dd u
       &=1+\frac43\log2+12\zeta'(-1),\label{pls:eq:E2m2}\\
 \int_0^1(u^2-u)\Phi(-1,3,1+u)\dd u
       &=\log(\pi/2)-\frac12,\label{pls:eq:E2m3}\\
 \int_0^1(u^2-u)\Phi(-1,4,1+u)\dd u
       &=\frac23\log2-\frac12.\label{pls:eq:E2m4}
\end{align}
For the first two, use \eqref{pls:eq:resvalue},
$\zeta(-1)=-1/12$, $\zeta(0)=-1/2$, and
$\zeta'(0)=-\tfrac12\log(2\pi)$. The last is
\eqref{pls:eq:eta}. Formula~\eqref{pls:eq:polygamma} turns all three
into explicit polygamma-difference integrals.

\section{Bernoulli selectors and generalized Stieltjes integrals}
\label{pls:sec:stieltjes}

Fix $N\ge2$, $a>0$, and
\begin{equation}\label{pls:eq:IN}
 I_N(s;a)=\int_0^1\bigl(B_N(u)-B_N\bigr)\zeta(s,a+u)\dd u.
\end{equation}
For odd $N\ge3$ the subtracted constant is zero. In all cases the
weight has mean $-B_N$. Write
\begin{equation}\label{pls:eq:bNgeneral}
 b_j=\frac{N!}{(N-j)!}B_{N-j}(1),\qquad 1\le j\le N.
\end{equation}
The resonant endpoint calculation gives
\begin{equation}\label{pls:eq:INclosed}
 \boxed{\displaystyle
 I_N(s;a)=-\frac{N!}{F_{N-1}(s)}\zeta(s-N,a)
          +\sum_{j=1}^N\frac{b_j a^{j+1-s}}{F_j(s)}.}
\end{equation}
This general-shift, unit-interval form has the same Bernoulli mechanism
as the predecessor's $[1,3]$ theorem \cite{Affine}. It is included to
make the confluence and primitive formulas below directly usable.

We use the Stieltjes convention
\begin{equation}\label{pls:eq:Stieltjes}
 \zeta(1+\varepsilon,x)=\frac1\varepsilon
     +\sum_{m=0}^\infty\frac{(-1)^m}{m!}\gamma_m(x)
                                               \varepsilon^m,
 \qquad \gamma_0(x)=-\psi(x).
\end{equation}
For $h\ge0$, define
\begin{equation}\label{pls:eq:Eh}
 E_h(\varepsilon)=\prod_{k=1}^h(1-\varepsilon/k)^{-1}
   =\exp\left(\sum_{r\ge1}\frac{H_h^{(r)}}r\varepsilon^r\right).
\end{equation}
Put
\begin{align}
 Q_N(\varepsilon;a)
 ={}&(-1)^N N E_{N-1}(\varepsilon)\zeta(1-N+\varepsilon,a)
        \notag\\
 &+\sum_{j=1}^N\frac{(-1)^j b_j a^j}{j!}
            E_j(\varepsilon)e^{-\varepsilon\log a}.
\label{pls:eq:QN}
\end{align}

\begin{proposition}[All weighted Stieltjes orders at one center]
\label{pls:prop:Stieltjes}
One has $Q_N(0;a)=-B_N$ and, for every $m\ge0$,
\begin{equation}\label{pls:eq:Stieltjesmoments}
 \boxed{\displaystyle
 \int_0^1\bigl(B_N(u)-B_N\bigr)\gamma_m(a+u)\dd u
   =(-1)^m m![\varepsilon^{m+1}]Q_N(\varepsilon;a).}
\end{equation}
This is a finite formula in $\zeta^{(k)}(1-N,a)$ for
$0\le k\le m+1$, powers of $a$ and $\log a$, rational Bernoulli
numbers, and generalized harmonic numbers.
\end{proposition}
\begin{proof}
At $s=1+\varepsilon$,
$F_j(s)^{-1}=(-1)^j E_j(\varepsilon)/(j!\varepsilon)$.
Substituting into \eqref{pls:eq:INclosed} gives exactly
$I_N(1+\varepsilon;a)=Q_N(\varepsilon;a)/\varepsilon$.
The compact integral has residue $\int_0^1(B_N(u)-B_N)\dd u=-B_N$.
Comparison with \eqref{pls:eq:Stieltjes}, justified uniformly on the
compact positive interval, proves both assertions.
\end{proof}

\subsection{Three explicit quartic evaluations}

For $N=4$ the weight is $u^2(1-u)^2$ and $b_2=2$, $b_3=12$,
$b_4=24$. In particular,
\begin{equation}\label{pls:eq:I4}
 I_4(s;a)=-\frac{24\zeta(s-4,a)}{F_3(s)}
       +\frac{2a^{3-s}}{F_2(s)}+\frac{12a^{4-s}}{F_3(s)}
                                      +\frac{24a^{5-s}}{F_4(s)}.
\end{equation}
At $a=1$ the coefficient generator simplifies to
\[
 Q_4(\varepsilon;1)
   =4E_3(\varepsilon)\zeta(-3+\varepsilon)
               +E_2(\varepsilon)-2E_3(\varepsilon)+E_4(\varepsilon).
\]
Expanding through degree two, using $\zeta(-3)=1/120$, yields
\begin{align}
 \int_0^1u^2(1-u)^2\gamma_0(1+u)\dd u
   &=4\zeta'(-3)-\frac1{45},\label{pls:eq:gamma0quartic}\\
 \int_0^1u^2(1-u)^2\gamma_1(1+u)\dd u
   &=-2\zeta''(-3)-\frac{22}{3}\zeta'(-3)+\frac5{432}.
                                            \label{pls:eq:gamma1quartic}
\end{align}
There are no unspecified Stieltjes constants on the right. A different
integer resonance of the same complex-order identity gives
\begin{equation}\label{pls:eq:trigammaquartic}
 \boxed{\displaystyle
 \int_0^1u^2(1-u)^2\psi'(1+u)\dd u
          =\frac{3\zeta(3)}{\pi^2}-\frac13.}
\end{equation}
Indeed, at $s=2+\varepsilon$ the constant term of
\eqref{pls:eq:I4} is $-1/3-12\zeta'(-2)$, and
$\zeta'(-2)=-\zeta(3)/(4\pi^2)$. Every integral here is ordinary
and absolutely convergent on a compact interval; no finite-part
convention is involved.

\section{Joint Stieltjes confluence at the untwisted point}
\label{pls:sec:confluence}

The change of polynomial basis at $z^L=1$ is algebraic. At $z=1$
there is a second, analytic issue: the Lerch function, entire in its
order before the limit, approaches a Hurwitz function with a pole.
The correct limiting object contains an explicit logarithmic term.
This section proves a joint holomorphic statement that controls every
order derivative at once.

Let $0<A<B$, and let $w\in L^1([A,B])$ be a fixed complex weight.
For polynomial weights, all the finite integrals appearing below can
in turn be reduced by endpoint formulas. Set
\begin{align}
 M_w(\mu)&=\int_A^B w(x)e^{\mu x}\dd x,\label{pls:eq:Mw}\\
 T_w(\mu,\varepsilon)&=\int_A^B
           w(x)\Phi(e^{-\mu},1+\varepsilon,x)\dd x,
              \qquad \mu>0.\label{pls:eq:Twmu}
\end{align}
For $\varepsilon$ near zero define
\begin{align}
 C(\varepsilon,\mu)
 &=\Gamma(-\varepsilon)\mu^\varepsilon+\frac1\varepsilon
                    \label{pls:eq:compensator}\\
 &=\frac{1-\exp\left((\log\mu+\gamma)\varepsilon
                 +\sum_{r\ge2}\zeta(r)\varepsilon^r/r\right)}
                 {\varepsilon}.
                    \label{pls:eq:Cexp}
\end{align}
Here $\gamma$ is Euler's constant, $\mu^\varepsilon$ uses the real
logarithm for $\mu>0$, and the singularity at $\varepsilon=0$ in
each display is removable. The second formula follows from the
classical Taylor expansion of $\log\Gamma(1-\varepsilon)$.

\begin{theorem}[Joint holomorphic confluence and every mixed jet]
\label{pls:thm:confluence}
The compensated transform
\begin{equation}\label{pls:eq:Rwdef}
 \mathscr R_w(\mu,\varepsilon)
     =T_w(\mu,\varepsilon)-M_w(\mu)C(\varepsilon,\mu)
\end{equation}
extends holomorphically to
$|\mu|<2\pi$, $|\varepsilon|<1$. Its exact expansion there is
\begin{align}
 \mathscr R_w(\mu,\varepsilon)
 ={}&\int_A^B w(x)e^{\mu x}
               \left(\zeta(1+\varepsilon,x)-\frac1\varepsilon\right)
                                                       \dd x\notag\\
 &+\sum_{k=1}^\infty\frac{(-\mu)^k}{k!}
       \int_A^B w(x)e^{\mu x}\zeta(1+\varepsilon-k,x)\dd x.
                                      \label{pls:eq:Rwseries}
\end{align}
The first bracket is interpreted by its analytic extension at zero.
For all nonnegative integers $h,m$,
\begin{align}
 \left.\partial_\mu^h\partial_\varepsilon^m
                      \mathscr R_w\right|_{(0,0)}
 ={}&(-1)^m\int_A^B w(x)x^h\gamma_m(x)\dd x\notag\\
 &+\sum_{k=1}^h(-1)^k\binom hk
     \int_A^B w(x)x^{h-k}\zeta^{(m)}(1-k,x)\dd x.
                                      \label{pls:eq:mixedjets}
\end{align}
The superscript on $\zeta$ differentiates its first argument.
\end{theorem}
\begin{proof}
The classical Lerch expansion \cite[Eq.~25.14.3\_1]{DLMFLerch},
with $z=e^{-\mu}$, is
\begin{equation}\label{pls:eq:classicalLerch}
 \Phi(e^{-\mu},1+\varepsilon,x)
 =e^{\mu x}\left\{\Gamma(-\varepsilon)\mu^\varepsilon
       +\sum_{k=0}^\infty\zeta(1+\varepsilon-k,x)
                                     \frac{(-\mu)^k}{k!}\right\}.
\end{equation}
Initially $0<\mu<2\pi$ and $\varepsilon$ avoids zero and the
other exceptional integers. Subtract $e^{\mu x}C$ and separate
$k=0$. The result is the integrand version of
\eqref{pls:eq:Rwseries}.

We record the normal-convergence justification needed for the joint
claim. Fix $0<\delta<1$ and a compact subdisk $|\mu|\le\rho<2\pi$.
For $k\ge2$, the Hurwitz Fourier formula at
$s=1+\varepsilon-k$, together with integer translation in the
positive argument, gives uniformly for $x\in[A,B]$ and
$|\varepsilon|\le\delta$
\begin{equation}\label{pls:eq:normalestimate}
 |\zeta(1+\varepsilon-k,x)|
 \le C\,\Gamma(k+\delta)(2\pi)^{-k}+C D^k
\end{equation}
for fixed $C,D$. To see the uniformity, reduce the argument to
$y\in[0,1]$. The Fourier sums are bounded by
$\sum_{n\ge1}n^{-k+\delta}$; their trigonometric factors are
bounded, and the gamma factor is controlled on the fixed
$\varepsilon$ disk. At $y=0$ use the continuous negative-order
limit. Translation back adds finitely many powers
$(y+j)^{k-1-\varepsilon}$, bounded by $C D^k$ since
$k-1-\delta>0$. The finite cases and the values $k=0,1$ are
handled separately after removing the single pole of the $k=0$ term.
The usual gamma-ratio estimate, or its elementary product form on
this fixed disk, now makes \eqref{pls:eq:normalestimate} divided by
$k!$ a polynomial in $k$ times $(2\pi)^{-k}$, plus an entire-series
term $C D^k/k!$.

Consequently the series in \eqref{pls:eq:Rwseries} is normally
convergent on compact sub-bidisks and uniformly in $x\in[A,B]$.
The factor $e^{\mu x}$ is bounded there. Integration against an
$L^1$ weight preserves normal convergence, so the right side is
jointly holomorphic. It agrees with \eqref{pls:eq:Rwdef} in the
initial domain and supplies the claimed extension. Cauchy's theorem
on smaller disks justifies every spectral derivative as well.

Finally, the coefficient of $\mu^h$ in
$e^{\mu x}(-\mu)^k/k!$ is
$(-1)^k x^{h-k}/(k!(h-k)!)$ when $k\le h$, and zero otherwise.
The regular Hurwitz expansion \eqref{pls:eq:Stieltjes} gives the
$k=0$ derivative. Multiplying by $h!$ proves
\eqref{pls:eq:mixedjets}.
\end{proof}

The estimate inside the proof is only used to justify an exact
identity and its differentiated limits; no optimization of bounds is
part of the result.

\subsection{Explicit logarithmic subtraction at every spectral order}

Define
\begin{equation}\label{pls:eq:Cm}
 C_m(\mu)=-m![t^{m+1}]
   \exp\left((\log\mu+\gamma)t+\sum_{r\ge2}\zeta(r)t^r/r\right).
\end{equation}
Then $C_m=\partial_\varepsilon^mC(0,\mu)$. With
$\ell=\log\mu+\gamma$, the first terms are
\begin{align}
 C_0(\mu)&=-\ell,\notag\\
 C_1(\mu)&=-\frac{\ell^2+\zeta(2)}2,\label{pls:eq:Cexamples}\\
 C_2(\mu)&=-\frac{\ell^3}3-\zeta(2)\ell-\frac{2\zeta(3)}3.
                    \notag
\end{align}
The exact all-order limit is therefore
\begin{equation}\label{pls:eq:limitjets}
 \boxed{\displaystyle
 \lim_{\mu\downarrow0}
 \left\{\int_A^B w(x)
        \left.\partial_s^m\Phi(e^{-\mu},s,x)\right|_{s=1}\dd x
                   -M_w(\mu)C_m(\mu)\right\}
   =(-1)^m\int_A^B w(x)\gamma_m(x)\dd x.}
\end{equation}
In particular, zeroth order requires addition of
$M_w(\mu)(\log\mu+\gamma)$, not subtraction of that expression.
The signs are fixed by the Laurent convention
\eqref{pls:eq:Stieltjes}.

\begin{corollary}[Mean-zero weights need no logarithmic subtraction]
\label{pls:cor:meanzero}
If $\int_A^B w(x)\dd x=0$, then the subtraction term in
\eqref{pls:eq:limitjets} tends to zero for every fixed $m$, and the
uncompensated limit exists with the displayed value. If that mean is
nonzero, the uncompensated zeroth-order integral diverges logarithmically.
\end{corollary}
\begin{proof}
In the mean-zero case $M_w(\mu)=O(\mu)$, whereas $C_m(\mu)$ is
a polynomial of degree $m+1$ in $\log\mu$. Their product tends
to zero. In the other case $M_w(\mu)$ tends to a nonzero constant
and $C_0=-\log\mu-\gamma$ is unbounded; the remainder has a finite
limit by the theorem.
\end{proof}

For $W_{2r}(u)=B_{2r}(u)-B_{2r}$, the mean is $-B_{2r}\ne0$.
Thus the endpoint-flat resonant weights in
Theorem~\ref{pls:thm:flatness} genuinely require the compensation.
Odd Bernoulli weights $B_{2r+1}(u)$ with $r\ge1$ have mean zero,
and their fixed-order Stieltjes limits need none. This is a distinction
between endpoint data and mean data, not a contradiction.

\subsection{A finite coefficient check at the origin}

For $w=1$ on $[1,2]$, \eqref{pls:eq:d0res} gives the original
Lerch integral through one polylogarithm derivative. The mass is
$M_w(\mu)=(e^{2\mu}-e^\mu)/\mu$. The extension satisfies
\begin{equation}\label{pls:eq:originexample}
 \mathscr R_w(0,0)=0,\qquad
 \partial_\mu\mathscr R_w(0,0)=\frac12\log(2\pi).
\end{equation}
Indeed the first value is
$-\int_1^2\psi(x)\dd x=0$. The second, by
\eqref{pls:eq:mixedjets}, is
$\int_1^2(-x\psi(x)-\zeta(0,x))\dd x$.
Integration by parts and $\zeta(0,x)=1/2-x$ reduce it to
$1+\int_1^2\log\Gamma(x)\dd x=\tfrac12\log(2\pi)$,
using the classical Raabe integral. This verifies both the twist
coordinate and the sign of the compensator.

\section{Anchored primitives in the Hurwitz argument}
\label{pls:sec:primitives}

Endpoint closure also survives repeated integration in the shift.
For this section $P$ and $L$ are fixed; abbreviate
$T(s;a)=\T_P(z,s;a,L)$. Let $a,a_0>0$, and let $q\ge1$.
The Pochhammer symbol in the next display is rising:
$(s-q)_q=(s-q)(s-q+1)\cdots(s-1)$.

\begin{theorem}[All anchored shift primitives]
\label{pls:thm:primitives}
Define meromorphically in $s$
\begin{equation}\label{pls:eq:Jq}
 J_q(s;a,a_0)=\frac{(-1)^q}{(s-q)_q}
 \left\{T(s-q;a)-\sum_{j=0}^{q-1}\frac{(a-a_0)^j}{j!}
                              \partial_a^jT(s-q;a_0)\right\}.
\end{equation}
Then
\begin{equation}\label{pls:eq:Janchor}
 \partial_a^qJ_q(s;a,a_0)=T(s;a),\qquad
 \left.\partial_a^jJ_q(s;a,a_0)\right|_{a=a_0}=0
                              \quad(0\le j<q).
\end{equation}
For $z\ne1$ all apparent spectral singularities of $J_q$ are
removable. For $z=1$ the only possible pole is at $s=1$, with
\begin{equation}\label{pls:eq:Jresidue}
 \Res_{s=1}J_q(s;a,a_0)
     =\frac{(a-a_0)^q}{q!}\int_0^L P(u)\dd u.
\end{equation}
Consequently, if $M=\int_0^LP(u)\dd u$, the coefficient
\begin{equation}\label{pls:eq:Stieltjesprimitive}
 \mathcal A_{m,q}(a,a_0)=(-1)^m m![\varepsilon^m]
 \left\{J_q(1+\varepsilon;a,a_0)
                  -\frac{M(a-a_0)^q}{q!\varepsilon}\right\}
\end{equation}
is the unique $q$th anchored primitive of
$\int_0^LP(u)\gamma_m(a+u)\dd u$.
\end{theorem}
\begin{proof}
Repeated argument differentiation gives
$\partial_a^qT(s-q;a)=(-1)^q(s-q)_qT(s;a)$.
The Taylor subtraction proves the anchors, and proves
\eqref{pls:eq:Janchor} away from the exceptional orders.
Alternatively the resulting function is the ordinary iterated integral
\begin{equation}\label{pls:eq:Jintegral}
 J_q(s;a,a_0)=\frac1{(q-1)!}\int_{a_0}^{a}
                           (a-v)^{q-1}T(s;v)\dd v,
\end{equation}
with the standard oriented integral if $a<a_0$. It gives analytic
continuation in $s$ and uniqueness with the anchors immediately.

The cancellation is also visible in the finite formula. For $z\ne1$,
$\Phi(z,-k,a)$ is a polynomial in $a$ of degree at most $k$:
expand $(n+a)^k$ in the convergent series when $|z|<1$, and then
continue to the unit circle away from one. At each zero
$s=1,\ldots,q$ of the Pochhammer factor, the Taylor remainder
therefore vanishes. For $z=1$, $\zeta(-k,a)=-B_{k+1}(a)/(k+1)$.
The zeros $s=2,\ldots,q$ still cancel; at $s=1$ the degree-$q$
term gives precisely \eqref{pls:eq:Jresidue}. The potential pole
at $s=q+1$ of the shifted $T(s-q)$ has residue independent of $a$,
so the zeroth Taylor term cancels it. This accounts for every apparent
singularity of \eqref{pls:eq:Jq}.

Finally apply \eqref{pls:eq:Stieltjes} to the ordinary iterated
integral \eqref{pls:eq:Jintegral}. Coefficient extraction gives
\eqref{pls:eq:Stieltjesprimitive}, including all anchors.
\end{proof}

For $q=1$ the Stieltjes formula is particularly transparent. With
$z=1$ and $T(s;a)=\int_0^LP(u)\zeta(s,a+u)\dd u$,
\begin{equation}\label{pls:eq:primitivefirst}
 \mathcal A_{m,1}(a,a_0)=\frac{(-1)^{m+1}}{m+1}
       \left\{\partial_s^{m+1}T(0;a)
                    -\partial_s^{m+1}T(0;a_0)\right\}.
\end{equation}
At $m=0$, Lerch's formula for $\zeta'(0,x)$ reduces it to
\begin{equation}\label{pls:eq:primitiveGamma}
 \mathcal A_{0,1}(a,a_0)
 =-\int_0^L P(u)
              \log\frac{\Gamma(a+u)}{\Gamma(a_0+u)}\dd u.
\end{equation}
At $m=1$ it is one half of the corresponding difference of weighted
$\zeta''(0,x)$ integrals. For selector weights, the endpoint identities
express these further in finitely many negative-order zeta derivatives.
This is an exact anchored calculus; an unspecified additive polynomial
would lose the normalization required by the source programme.

\section{A related open coefficient problem and normalization corrections}
\label{pls:sec:coefficientproblem}

The literature search located a concrete complementary question in
Pokaz, Ribi\v ci\'c Penava, and Simsek, arXiv:2609.13796v1
\cite{PPS}. On page~9, the authors ask for the coefficients in a
finite Hurwitz--Lerch reduction of a nested series. Their examples
cover the first few values of a discrete parameter. The complete
answer is a binomial row. This is an elementary partial-fraction
resolution of the question as posed in that version, not a claim that
the partial-fraction method itself is new.

Put $n=v-1\ge1$, $a=x-1>0$, and define
\begin{equation}\label{pls:eq:hmnodes}
 h_m^{[n]}=\sum_{m_1+\cdots+m_n=m}\prod_{j=1}^n j^{-m_j},
 \qquad
 Z_v(s;a)=-\frac{(s)_v}{(v-1)!}
              \sum_{m\ge0}\frac{h_m^{[v-1]}}{(a+m)^{s+v}}.
\end{equation}
Here $(s)_v$ is rising and $\Re(s+v)>1$ initially. The change
$m_1=k_1$, $m_j=k_j-k_{j-1}$, and
$m_{v-1}=m-k_{v-2}$ identifies this with Definition~1 of
\cite{PPS}, with $Z_v(s;a)=Z(s;a+1,v)$ in that paper's notation.

\begin{theorem}[The coefficient question in \cite{PPS}]
\label{pls:thm:alphas}
For every $v\ge2$ and $1\le j\le v-1$, set
\begin{equation}\label{pls:eq:alphas}
 \boxed{\alpha_j=(-1)^{j-1}\binom{v-1}{j}.}
\end{equation}
Then, identically in $s$ by continuation,
\begin{equation}\label{pls:eq:Zclosed}
 \boxed{\displaystyle
 Z_v(s;a)=-\frac{(s)_v}{(v-1)!}
           \sum_{j=1}^{v-1}(-1)^{j-1}\binom{v-1}{j}
                              \Phi(1/j,s+v,a).}
\end{equation}
The $j=1$ term is $(v-1)\zeta(s+v,a)$, as required in the
source question. The function $Z_v(s;a)$ is entire in $s$.
\end{theorem}
\begin{proof}
The generating function for the finite-composition coefficients is
\[
 \sum_{m\ge0}h_m^{[n]}t^m=\prod_{j=1}^n(1-t/j)^{-1}.
\]
The coefficient of $(1-t/j)^{-1}$ in its partial-fraction expansion is
\[
 \prod_{\substack{1\le k\le n\\k\ne j}}(1-j/k)^{-1}
   =\frac{n!/j}{(-1)^{j-1}(j-1)!(n-j)!}
   =(-1)^{j-1}\binom nj.
\]
Both sides have the same simple poles and principal parts and tend
to zero at infinity, so the partial-fraction identity is exact.
Comparing coefficients gives
\begin{equation}\label{pls:eq:hmclosed}
 h_m^{[n]}=\sum_{j=1}^n(-1)^{j-1}\binom nj j^{-m}.
\end{equation}
These coefficients are bounded as $m\to\infty$; hence the defining
Dirichlet series is absolutely convergent for $\Re(s+v)>1$.
Substitute \eqref{pls:eq:hmclosed} and interchange the finite sum to
obtain \eqref{pls:eq:Zclosed}. Every term with $j\ge2$ is entire.
The sole Hurwitz pole at $s=1-v$ is canceled by the simple zero of
$(s)_v$. Thus $Z_v$ is entire.
\end{proof}

\subsection{Every jet at the canceled Hurwitz pole}

The same harmonic factor as in the selector calculus appears at
$s_0=1-v$. Write $n=v-1$ and
$D_n(\varepsilon)=\prod_{k=1}^n(1-\varepsilon/k)=E_n(\varepsilon)^{-1}$.
Then \eqref{pls:eq:Zclosed} is equivalent near $s_0$ to
\begin{align}
 Z_v(1-v+\varepsilon;a)
 ={}&(-1)^vD_{v-1}(\varepsilon)
 \Bigg[(v-1)\notag\\[-2pt]
 &+\varepsilon\left\{(v-1)\sum_{m\ge0}
          \frac{(-1)^m\gamma_m(a)}{m!}\varepsilon^m
       +\sum_{j=2}^{v-1}\alpha_j\Phi(1/j,1+\varepsilon,a)
                   \right\}\Bigg].
\label{pls:eq:Zjets}
\end{align}
This convergent local identity proves every spectral derivative by
finite coefficient extraction. In particular,
\begin{align}
 Z_v(1-v;a)&=(-1)^v(v-1),\label{pls:eq:Zvalue}\\
 Z_v'(1-v;a)&=(-1)^v\left\{(v-1)(\gamma_0(a)-H_{v-1})
                 +\sum_{j=2}^{v-1}\alpha_j\Phi(1/j,1,a)\right\}.
                                      \label{pls:eq:Zderivative}
\end{align}
For example $Z_3'(-2;1)=3-2\gamma+2\log2$. Thus the coefficient
resolution also supplies exact Stieltjes and logarithmic evaluations,
not just an interpolation of the first few coefficient rows.

\subsection{Normalization fixed by the defining series}

The definition \eqref{pls:eq:hmnodes}, not an inconsistent later
normalization, is used throughout. Direct Mellin integration of that
definition gives, for $a>0$ and $\Re(s+v)>1$,
\begin{equation}\label{pls:eq:ZcorrectMellin}
 \boxed{\displaystyle
 Z_v(s;a)=-\frac{1}{(v-1)!\,\Gamma(s)}
       \int_0^\infty
       \frac{t^{s+v-1}e^{-at}}{\prod_{j=1}^{v-1}(1-e^{-t}/j)}\dd t.}
\end{equation}
At poles of $\Gamma(s)$ this is interpreted by its entire reciprocal.
To prove it, expand each geometric denominator, integrate
$t^{s+v-1}e^{-(a+m)t}$, and use
$(s)_v=\Gamma(s+v)/\Gamma(s)$ by continuation.

In the inspected version of \cite{PPS}, equation~(17) on page~11
has the prefactor $-(v-1)!/\Gamma(s)$ instead. Relative to its own
Definition~1 this differs by $((v-1)!)^2$. The subsequent $v=2$
integral on that page also has an incorrect sign and gamma factor.
Its corrected form is
\begin{equation}\label{pls:eq:correctv2}
 \int_0^\infty\frac{t^{s+1}e^{-at}}{1-e^{-t}}\dd t
             =\Gamma(s+2)\zeta(s+2,a),
 \qquad a>0,\quad\Re s>-1.
\end{equation}
This follows termwise from a positive kernel when $s$ is real; a
negative right side in that domain cannot be correct. Both source
pages were checked in the rendered PDF, not only in extracted text.

One can also settle the sign of the polynomial interpolation without
a contour orientation convention. Retain the source generating
function
\begin{equation}\label{pls:eq:ellgen}
 F(t,x;v)=\frac{t^v e^{xt}}{\prod_{j=0}^{v-1}(e^t-j)}
           =\sum_{k\ge0}\ell_k(x;v)\frac{t^k}{k!}.
\end{equation}
Algebra gives
\[
 F(-t,a+1;v)=
 -\frac{t^v e^{-at}}{(v-1)!\prod_{j=1}^{v-1}(1-e^{-t}/j)}.
\]
Therefore \eqref{pls:eq:ZcorrectMellin} is
$\Gamma(s)^{-1}\int_0^\infty t^{s-1}F(-t,a+1;v)\dd t$.
Subtracting a finite Taylor polynomial at zero continues this Mellin
integral. Its residue at $s=-k$ is
$(-1)^k\ell_k(a+1;v)/k!$, whereas
$\Gamma(s)^{-1}\sim(-1)^k k!(s+k)$. Hence
\begin{equation}\label{pls:eq:correctinterpolation}
 \boxed{Z_v(-k;a)=\ell_k(a+1;v),\qquad k\ge0.}
\end{equation}
The minus sign in equation~(19) of \cite{PPS} is incompatible with
Definition~1. Already $v=2$, $k=1$ gives $Z_2(-1;a)=1$ and
$\ell_1(a+1;2)=1$. This correction is scoped to that source and
normalization; it is not an allegation about the ProveIt manuscript.

\section{Source audit, verification, and integration status}
\label{pls:sec:audit}

\subsection{What was inspected}

The canonical manuscript README and the incoming-directory metadata were
read through the connected GitHub tools on 11 October 2026. Selected
harmonic-tail sections were read at commit
\texttt{47f2d0da4f1c0e19f9ec8aecb39112756f167d82}.
A later repository search used a newer indexed commit. These observations
are not represented as a single frozen snapshot of a changing
\src{main} branch.

The precise mathematical predecessor for this article is the 1,702-line
standalone TeX source of \emph{Affine Gamma Polarization}, obtained from
the user's Library. Its sections on Lerch endpoints, Bernoulli selectors,
Stieltjes coefficients, and further questions were inspected. Its SHA-256
and Library version are recorded in \src{source_manifest.json}. The
GitHub incoming directory independently lists the corresponding named
archive, but its binary contents were not extracted in this run. We do
not assert a byte-for-byte comparison between the Library TeX and every
member of that archive. No external source PDF or predecessor TeX is
redistributed in this package.

For the external coefficient question, the defining series, the open
question, the Mellin normalization, and the interpolation statement in
arXiv:2609.13796v1 were checked in both the HTML and the rendered PDF.
The relevant printed pages are 3, 9, 11, and 12. The corrections in
Section~\ref{pls:sec:coefficientproblem} are tied to that version's
Definition~1. They do not silently change its definition to match a
later, inconsistent prefactor.

\subsection{Corrections and limitations}

There is no false theorem identified in the reviewed selector portion
of the incoming ProveIt report. Its open question genuinely leaves room
for the classification proved here. Three cautions are nevertheless
necessary when integrating the continuation.

First, one cannot deform a double-flat Bernoulli selector into a
nonresonant double-flat single-order polynomial. The obstruction is
Theorem~\ref{pls:thm:flatness}, not a failure to find the appropriate
normalization. The exact replacement is the finite mixture
\eqref{pls:eq:mixing}. Second, a fixed-degree Frobenius--Euler selector
cannot have a nonconstant finite scalar limit at resonance; the
coefficientwise result in Section~\ref{pls:sec:mixing} explains why
singular coefficients in the mixture must cancel collectively. Third,
the limit $z\to1$ cannot be substituted into a spectral jet at $s=1$
without its pole/contact subtraction. Equations
\eqref{pls:eq:compensator}--\eqref{pls:eq:limitjets} supply the complete
subtraction, including Euler's constant and all higher zeta terms.

The external corrections are independent: relative to the defining
series, equation~(17) of \cite{PPS} has an erroneous factorial factor;
its displayed $v=2$ integral has the wrong sign and gamma factor; and
equation~(19) has the wrong interpolation sign. The replacement formulas
\eqref{pls:eq:ZcorrectMellin}, \eqref{pls:eq:correctv2}, and
\eqref{pls:eq:correctinterpolation} are derived from the definition and
checked independently. This is a scoped audit, not a claim to have
verified every formula in that paper or in the ProveIt volumes.

\subsection{Proof status and dependence}

\begin{center}
\small
\begin{tabular}{@{}p{0.30\textwidth}p{0.61\textwidth}@{}}
\toprule
Result & Scope and principal justification\\
\midrule
Selector classification & Every integer length, $a>0$, and
$0<|z|\le1$; endpoint identity and functional independence.\\
Endpoint obstruction & Every degree; Eulerian interlacing and the
classical Appell boundary relation.\\
Resonant support collapse & Exact identity at every even degree;
Frobenius--Euler endpoint reconstruction.\\
Spectral and Stieltjes jets & Every jet order; analytic removability,
finite coefficient extraction, and normal convergence.\\
Anchored primitives & Every primitive order; Taylor subtraction and
ordinary iterated integration.\\
External coefficient question & Every $v\ge2$; explicit rational
partial fractions, not extrapolation of finite rows.\\
Arithmetic independence & Not asserted after specializing $s$.
Functional independence over $\C(s)$ is a different theorem.\\
\bottomrule
\end{tabular}
\end{center}

The Frobenius--Euler basis theorem is already present in Kim and Kim
\cite{KimKim}; the endpoint transform uses classical Hurwitz/Lerch
calculus. Integrals of the Hurwitz zeta function, log-Gamma functions,
and negative-order polygamma functions have a substantial prior
literature, including Espinosa and Moll \cite{EM1,EM2}. The examples
here are consequences in a specified selector system, not claims that
weighted Gamma integrals or their general method were previously
unknown. Likewise, the coefficient solution in
Section~\ref{pls:sec:coefficientproblem} is a complete answer to a
stated question, but uses elementary and classical partial fractions.
Global bibliographical novelty remains unverified.

\subsection{Exact checks}

The scripts use SymPy 1.14.0 and rational arithmetic. The original
selector suite contains 667 assertions. The separately implemented
coefficient-problem suite contributes 302 more, for a total of
\textbf{969 finite exact assertions}. Six of these are mutation controls
which deliberately insert an incorrect expression and verify its
rejection. The checks include the following independent constructions.

\begin{center}
\small
\begin{tabular}{@{}p{0.70\textwidth}r@{}}
\toprule
Selector-suite family & Assertions\\
\midrule
Frobenius--Euler endpoint derivatives & 180\\
Appell derivatives and endpoint reconstruction & 74\\
Bernoulli resonant derivatives & 90\\
Eulerian endpoint identities, factorial values, consecutive gcds & 29\\
Algebraic twist and exact support strata & 16\\
Finite-mixture identities and endpoint multiplicities & 25\\
Removable-integer numerator cancellation & 135\\
Generalized-harmonic coefficient recurrence & 112\\
Printed Stieltjes/trigamma identities & 4\\
Mutation controls & 2\\
\midrule
Selector subtotal & 667\\
\midrule
External rational partial fractions & 10\\
Nested-composition coefficients versus binomial rows & 210\\
Defining generating series versus continued integer values & 78\\
External normalization/sign mutation controls & 4\\
\midrule
Combined total & 969\\
\bottomrule
\end{tabular}
\end{center}

The interpolation tests compute the coefficients of the source
function $F(t,2;v)$ by rational power-series convolution. The other
side is computed from continued zeta/Lerch values, with the canceled
pole handled separately. Thus the sign tests do not simply evaluate
two copies of the same Mellin formula. The external coefficient tests
also generate the nested coefficients by multiplying geometric series
rather than by using their binomial answer.

\subsection{Numerical diagnostics}

The two diagnostic suites use mpmath 1.3.0 at 75 decimal digits of
working precision. There are \textbf{33 comparisons}: 27 for the
selector calculus and six for the external coefficient and Mellin
identities. The largest recorded scaled discrepancy is approximately
$3.075\times10^{-36}$. This comes from the independent quartic
first-Stieltjes integral, computed with a finite Gauss--Legendre rule;
its acceptance tolerance is $10^{-32}$. Individual absolute and scaled
discrepancies and all tolerances are recorded in the result JSON files.
Working precision is not an accuracy certificate.

The comparisons include direct geometric sums against endpoint
formulas, removable-order limits, an order derivative, ordinary
Euler-weighted digamma and log-Gamma quadratures, three quadratic
polygamma identities, the quartic Stieltjes and trigamma formulas,
compensated Lerch expansions, a mixed confluence derivative, and an
anchored primitive. The external Mellin comparisons use an independently
accumulated nested-coefficient series; its constant tail is summed by
Hurwitz zeta and its remaining geometric tail is truncated. The
canceled-pole jet is checked by Cauchy differentiation on a small
circle, avoiding direct multiplication of a rounded zero by a pole.

These are floating-point diagnostics, not interval-certified enclosures.
They can detect index, sign, and implementation errors. They do not
replace any all-order proof, prove numerical period independence, or
certify global novelty. No theorem prover has been used.

\subsection{Integration proposal}

Retain this article and its scripts as a separate continuation report,
with the ledger marking the integer-length version of the source's
questions~4 and~5 as resolved in the stated functional sense.
The rational-endpoint, noninteger-length question remains separate.
The primitive and joint-jet theorems advance question~6 but do not
supply its requested compact degree/order bivariate generating function
or a complete Barnes-Gamma classification.

Every label in this source uses the \src{pls:} prefix. The package
includes a short integration fragment, an explicit source manifest,
reproduction instructions, and exact/numerical receipts. No repository
files have been edited. In particular, this delivery does not change
the conjectural status of $S_6$, the revised $S_8$, or the general
unequal-tail optimal-prefix problem.

\section{Further research questions}
\label{pls:sec:questions}

\paragraph{1. Prescribed spectral support with higher endpoint vanishing.}
Classify polynomial weights selecting two or more prescribed shifted
orders while vanishing to a specified order at both endpoints.
Theorem~\ref{pls:thm:flatness} settles the one-coordinate obstruction.
For larger support, the exact problem is the rank of an endpoint matrix
of Frobenius--Euler polynomials. Determine its exceptional algebraic
twists and prove minimal support in a declared degree range. Counting
equations alone does not prove the necessary minors are nonzero.

\paragraph{2. Rational interval lengths and cyclotomic coordinates.}
For noninteger rational $L$, translation no longer returns the same
Hurwitz argument modulo a finite prefix. Determine a complete selector
basis with a finite vector of rational residue-class arguments, then
use roots-of-unity filters to convert it to polylogarithms or Dirichlet
$L$-functions. The coefficient field and the permitted finite remainder
must be fixed before stating uniqueness. Integer $L$ is completely
covered here; this different endpoint geometry is not.

\paragraph{3. Coalescing geometric poles.}
Replace the distinct nodes $1,1/2,\ldots,1/(v-1)$ in
Section~\ref{pls:sec:coefficientproblem} by arbitrary nodes with
multiplicities. Partial fractions predict finite derivatives in the
Lerch twist, while the collision limit is governed by divided
differences. Prove a stable all-jet confluence theorem and determine
which normalizations give particularly simple Gamma-product or
Stieltjes interpolation values. Keep the original defining series as
the normalization anchor.

\paragraph{4. Resonant roots of unity and simultaneous degenerations.}
At $z^L=1$ with $z\ne1$, spectral support collapses but there is no
Hurwitz pole. Combine the exact support law with finite Fourier
transforms to distinguish this pure endpoint resonance from the
additional logarithmic contact at $z=1$. Study simultaneous variation
of the interval shift, twist, and spectral order, with an explicit
finite subtraction at every singular face.

\paragraph{5. A compact degree--Stieltjes generating identity.}
The coefficient rules in Sections~\ref{pls:sec:jets} and
\ref{pls:sec:stieltjes} compute every fixed degree and every Stieltjes
order. Seek a bivariate generating identity that retains a single
spectral center and exposes the complete-harmonic coefficients without
a separate finite sum at each degree. Determine which anchored
primitives close in a specified finite family of multiple-Gamma or
zeta-derivative functions. This is the remaining substantive part of
the predecessor's question~6.

\paragraph{6. Nonlinear Gamma sectors with exact selectors.}
The incoming affine Gamma square-sector formula motivates the beta
and Bernoulli weights used here. Transport multi-coordinate selectors
back through that normal-form construction, retaining every endpoint
term. Seek complete cubic or quartic derivative-sector identities,
not only diagonal squares or low-degree examples. A useful outcome
would be a generating formula whose Hurwitz/Lerch transform has a
provably minimal, prescribed functional spectrum.

\paragraph{7. The remaining harmonic-tail and polylogarithm questions.}
The general unequal-even-tail optimal-prefix conjecture concerns
nonvanishing of specific rational coefficient determinants and is not
implied by the independence lemma here. Find an exact representation
of those determinants or a recurrence proving nonvanishing. Separately,
the $S_6$ and revised $S_8$ period candidates require exact parameter
identities with controlled boundary constants. A failed relation
search, or an extremely small residual, does not settle either problem.

\paragraph{8. Proof-carrying symbolic integration.}
Formalize the endpoint map, its basis inversion, the rational
partial-fraction theorem, and the finite harmonic coefficient
recurrence in a proof assistant. These finite algebraic cores can be
checked independently of the harder analytic continuation library.
Then add normal convergence, argument differentiation, and anchored
primitive uniqueness. The resulting compiler should emit both an
identity and a machine-checkable proof of its domain and normalization.

\subsection*{Conclusion}

Paired endpoint derivatives turn an integral-classification problem
into finite interpolation, while functional independence proves that
no hidden order has been discarded. Resonance changes the rank and
therefore the degree of the selector; it also explains why a finite
mixture, rather than scalar normalization, is needed to recover a
nonconstant Bernoulli weight. Complete-harmonic coefficients and a
single explicit logarithmic counterterm then connect this structure to
Gamma integrals, all Stieltjes jets, and anchored primitives. The same
attention to finite coefficients resolves the external nested-series
question and fixes its normalization. All claims retain their precise
algebraic, analytic, and computational scopes.

\begin{thebibliography}{99}
\raggedright
\bibitem{Manuscript}
ProveIt Contributors, \emph{Polylogarithms and their Arithmetic Bridges},
canonical manuscript README and selected harmonic-tail continuations.
Repository sources inspected 11 October 2026; moving-branch and pinned
reads are distinguished in Section~\ref{pls:sec:audit} and the source manifest.
\url{https://github.com/VladimirReshetnikov/ProveIt/tree/main/Analysis/Polylogarithms/docs/manuscript}.

\bibitem{Affine}
\emph{Affine Gamma Polarization}, ProveIt research continuation,
11 October 2026. Standalone TeX, Library version 1, 1,702 lines.
Sections 7, 9, and 11 are the principal predecessor material used here.
Named incoming archive:
\src{ProveIt_Affine_Gamma_Polarization_2026-10-11.zip}.
\url{https://github.com/VladimirReshetnikov/ProveIt/tree/main/docs/incoming}.
The standalone-source hash is recorded in \src{source_manifest.json}.

\bibitem{KimKim}
Dae San Kim and Taekyun Kim,
\emph{Some identities of Frobenius--Euler polynomials arising from
Frobenius--Euler basis}, arXiv:1211.6640 (2012), especially Theorem 1.
\url{https://arxiv.org/abs/1211.6640}.

\bibitem{Chow}
Chak-On Chow, \emph{New proofs of interlacing of zeros of Eulerian
polynomials}, Journal of Mathematical Analysis and Applications
\textbf{510} (2022), 126019.
\url{https://doi.org/10.1016/j.jmaa.2022.126019}.

\bibitem{DLMFLerch}
NIST Digital Library of Mathematical Functions,
Section 25.14, \emph{Lerch's Transcendent}, especially the series,
Mellin representation, translation, and logarithmic expansion.
\url{https://dlmf.nist.gov/25.14}. Consulted 11 October 2026.

\bibitem{DLMFHurwitz}
NIST Digital Library of Mathematical Functions,
Section 25.11, \emph{Hurwitz Zeta Function}, including argument
derivatives, special values, and Lerch's log-Gamma formula.
\url{https://dlmf.nist.gov/25.11}. Consulted 11 October 2026.

\bibitem{DLMFInteger}
NIST Digital Library of Mathematical Functions,
Section 25.6, \emph{Integer Arguments} of the Riemann zeta function.
\url{https://dlmf.nist.gov/25.6}. Consulted 11 October 2026.

\bibitem{DLMFPsi}
NIST Digital Library of Mathematical Functions,
Section 5.15, \emph{Polygamma Functions}.
\url{https://dlmf.nist.gov/5.15}. Consulted 11 October 2026.

\bibitem{EM1}
Olivier R. Espinosa and Victor H. Moll,
\emph{On some Definite Integrals involving the Hurwitz Zeta function},
arXiv:math/0012078 (2000).
\url{https://arxiv.org/abs/math/0012078}.

\bibitem{EM2}
Olivier R. Espinosa and Victor H. Moll,
\emph{On some integrals involving the Hurwitz zeta function: part 2},
arXiv:math/0107082 (2001).
\url{https://arxiv.org/abs/math/0107082}.

\bibitem{PPS}
Dora Pokaz, Mihaela Ribi\v ci\'c Penava, and Yilmaz Simsek,
\emph{Certain families of series involving Hurwitz and Lerch zeta functions},
arXiv:2609.13796v1, 12 September 2026.
The open coefficient question is on page 9; the corrected normalization
and interpolation statements are on pages 11--12. Version-specific
corrections are proved in Section~\ref{pls:sec:coefficientproblem}.
\url{https://arxiv.org/abs/2609.13796v1}.
\end{thebibliography}
\end{document}
